\documentclass{article}
\usepackage[utf8]{inputenc}
\usepackage[T1]{fontenc}
\usepackage{amsmath,amssymb}
\usepackage{booktabs}
\usepackage{graphicx}
\usepackage{xcolor}
\usepackage{natbib}
\usepackage[margin=1in]{geometry}
\usepackage[colorlinks=true,allcolors=blue!60!black]{hyperref}

% A visibly marked slot for a result that requires the LLM runs. Never silently
% empty: a placeholder that looks like prose is how unrun experiments end up
% being reported.
\newcommand{\pending}[1]{%
  \par\medskip\noindent\fcolorbox{red!60!black}{red!4}{%
    \parbox{0.97\linewidth}{\small\textbf{[PENDING LLM RUN]}~#1}}\par\medskip}

\title{\textsc{NudgeSim}: Costly Peer Correction as a Public Good\\
in Societies of LLM Agents}
\author{}
\date{}

\begin{document}
\maketitle

\begin{abstract}
Accurate information in a social network is a public good: verifying a claim is
privately costly while a less-polluted feed benefits everyone, so sharing
engaging but unverified content is free-riding and publicly correcting someone
else is costly peer sanctioning. That makes counter-misinformation an instance
of the second-order free-rider problem rather than a classification task. We
introduce \textsc{NudgeSim}, an environment in which a seven-agent society plays
this dilemma over real rumour-cascade topologies with human-vetted false claims,
under an explicit per-round payoff function, and in which an embedded
``Devil's Advocate'' intervener is manipulated along timing and tone. The
headline outcome is not what agents say they believe but what they pay for: the
rate at which citizens voluntarily absorb the correction cost once an intervener
is present.

Building the environment produced two findings that hold independently of which
model runs the agents, and that apply to any study in this line. First,
\emph{how a reply tree is turned into a visibility graph is an unstated
researcher degree of freedom that moves the headline effect by a factor of two
and a half}; treating the reply tree itself as the visibility graph---the
default---leaves no citizen able to observe another in 88.6\% of real
\textsc{Pheme} neighbourhoods, deleting the mechanism the study exists to
measure. Second, \emph{summing a payoff function is not a welfare measure in any
design with costly sanctioning}: total payoff ranks every working intervention
far below doing nothing, because the sanction costs it charges are precisely
what a working intervention generates. We report both with the environment, an
analysis pipeline validated against declared data-generating processes, and a
measured cost model for running the grid on language models.
\end{abstract}

%======================================================================
\section{Introduction}

Cooperation is expensive for the individual and valuable for the group, and
human societies sustain it in part through costly sanctioning: people pay out of
pocket to punish defectors \citep{fehr2002,guerek2006,ostrom1990}. A recent line
of work asks whether societies of large language model agents do the same,
placing them in common-pool resource games \citep{piatti2024govsim}, public
goods games with institutional choice \citep{piedrahita2025sanctsim}, morally
charged dilemmas \citep{backmann2025moralsim}, and mechanism-augmented matrix
games \citep{tewolde2026coopeval}. The consistent finding is that model
behaviour in these settings is neither uniform nor predictable from capability
alone---reasoning-tuned models free-ride where their non-reasoning counterparts
cooperate \citep{piedrahita2025sanctsim,tewolde2026coopeval}.

Misinformation is a social dilemma of exactly this shape, and it has not been
studied as one. Accurate information in a feed is a public good. Verification is
the contribution: privately costly, socially beneficial. Sharing engaging but
unverified content is free-riding---it collects the conformity reward
immediately at no verification cost while the veracity penalty is deferred and
uncertain. Publicly challenging another agent's claim is costly peer punishment,
which reintroduces the second-order dilemma: everyone wants the network
corrected, nobody wants to pay to correct it.

\textsc{NudgeSim} places an embedded AI intervener inside this dilemma and asks
a question the simulation literature on counter-misinformation has not: not
whether an intervention reduces propagation, but whether it \emph{catalyses or
crowds out} the society's own willingness to pay for correction. Crowding-out is
a documented effect of external enforcement on human intrinsic motivation
\citep{fehr2003,bowles2008}; whether it appears in an LLM society, and whether
the intervener's tone decides which way it goes, is the central question.

This paper reports the environment and the findings that building it produced.
Our contributions:

\begin{enumerate}
\item \textbf{\textsc{NudgeSim}}, an environment casting counter-misinformation
  as a public goods game with costly peer sanctioning, seeded with human-vetted
  false claims from \textsc{Liar} \citep{wang2017liar} and wired from the reply
  structure of real rumour cascades in \textsc{Pheme} \citep{kochkina2018pheme}
  (6{,}425 thread structures, 5{,}741 of which parse into usable reply trees), with an explicit payoff function and outcomes read
  directly off the action log rather than from a model judge
  (\S\ref{sec:env}).

\item \textbf{Reply trees are not visibility graphs, and the difference is the
  effect size.} Sampling a seven-node neighbourhood from a real \textsc{Pheme}
  cascade and treating reply edges as visibility edges leaves 88.6\% of
  neighbourhoods with no citizen able to see any other citizen. We measure three
  translations of the same corpus on the same seeds: the intervention's measured
  effect on peer correction runs $+0.038$, $+0.081$ and $+0.095$ depending only
  on which is chosen (\S\ref{sec:rq1}).

\item \textbf{Total payoff is not a welfare measure under costly sanctioning.}
  In \textsc{NudgeSim} the intervention cuts propagation and nearly quadruples
  peer correction, and $\sum\pi$ ranks it $122$ points below doing nothing.
  Removing sanction costs does not repair it, because the conformity term is
  maximised by the arm in which nobody disagrees. Only the veracity term tracks
  the good being provided (\S\ref{sec:rq2}).

\item \textbf{A validated analysis pipeline and a measured budget.} The analysis
  is run against three declared data-generating processes---effect present,
  inflated, and set to zero---and must recover them in that order
  (\S\ref{sec:sensitivity}). We meter the prompts the environment assembles and
  price the grid, finding that response caching, the lever such studies usually
  name, is worth nothing here, while the choice to enable extended thinking is
  worth an order of magnitude (\S\ref{sec:cost}).
\end{enumerate}

The model benchmark the environment is built for is not yet run; every slot it
fills is marked in the text and left empty rather than estimated.

%======================================================================
\section{Related Work}

\paragraph{LLM agents in social dilemmas.}
\textsc{GovSim} \citep{piatti2024govsim} places LLM agents around a regenerating
common-pool resource and finds that most models drive it to collapse, that
free-form communication is the largest single lever on survival, and that a
selfish newcomer inserted into a cooperating community degrades it.
\textsc{SanctSim} \citep{piedrahita2025sanctsim} adapts the institutional-choice
public goods game of \citet{guerek2006} and finds that reasoning models
free-ride where traditional models cooperate, and that LLMs reward where humans
punish. \textsc{MoralSim} \citep{backmann2025moralsim} varies the moral framing
of a fixed payoff structure and estimates per-factor average treatment effects.
\textsc{CoopEval} \citep{tewolde2026coopeval} benchmarks cooperation-sustaining
mechanisms across matrix games. \textsc{NudgeSim} is an instance of this
programme with information as the public good, and asks the second-order
question of an \emph{intervention} rather than of a population.

\paragraph{Counter-misinformation simulation.}
Agent-based studies of correction typically measure propagation or a
belief-like endorsement of the false claim, with no explicit payoff and no
equilibrium benchmark, so an intervention that reduces sharing counts as a
success regardless of what it does to the society's own willingness to correct.
\textsc{NudgeSim}'s headline outcome is that willingness.

\paragraph{Network data for agent societies.}
Studies in this area routinely seed agent adjacency from real cascade data.
\S\ref{sec:visibility} shows that the step from such data to an adjacency
matrix is neither neutral nor usually reported.

%======================================================================
\section{The \textsc{NudgeSim} Environment}
\label{sec:env}

\subsection{Economic background}

Let a society of $N$ agents face a single claim of unknown veracity over $T$
rounds. Each round every agent simultaneously chooses one action from
$\{\textsc{Share}, \textsc{Endorse}, \textsc{Challenge}, \textsc{Ignore}\}$
after observing what its neighbours did. \textsc{Share} and \textsc{Endorse}
propagate the claim; \textsc{Challenge} disputes it and costs the challenger.

\subsection{Environment dynamics}

Agent $i$'s payoff in round $t$ is
\begin{equation}
\pi_{i,t} \;=\; \beta \cdot C_{i,t} \;+\; \gamma \cdot V_i
           \;-\; \kappa \cdot \mathbf{1}[a_{i,t} = \textsc{Challenge}]
           \;-\; \delta \cdot D_{i,t},
\label{eq:payoff}
\end{equation}
where $C_{i,t}$ is the conformity reward paid immediately for agreeing with the
local majority, $V_i$ is the veracity reward resolved only at the horizon,
$\kappa$ is the verification and social-friction cost of challenging, and
$D_{i,t}$ is reputation damage taken by agents publicly backing a claim that
resolves false. The mapping is deliberate: $\beta C$ pays for engagement now,
$\gamma V$ pays for accuracy later, and $\kappa$ is what makes correction a
contribution nobody is individually motivated to make.

Three roles occupy the society. Five \emph{citizens} are payoff-responsive and
free. One \emph{disseminator} injects and re-frames the claim every round and is
by construction incorrigible. One \emph{intervener}---the Devil's
Advocate---is stipulated to absorb $\kappa$ whenever it challenges, which is
what makes it an exogenous sanctioner rather than a participant in the dilemma.
It first judges whether the claim is false and may decline to act, so that the
true-claim placebo arm tests specificity rather than assuming it.

\paragraph{Communication.}
The two reference designs answer the communication question oppositely:
\textsc{SanctSim} gives agents no natural-language channel at all, so every
behavioural difference is attributable to the payoff structure; \textsc{GovSim}
gives them free-form group chat and finds that removing it changes the headline
outcome by 21 points. \textsc{NudgeSim} splits them by channel. The
\emph{decision} channel is \textsc{SanctSim}-style: the payoff ledger in
Eq.~\ref{eq:payoff} is a function of actions alone, so nothing an agent writes
changes another's payoff except through the action it accompanies. The
\emph{observable} channel is \textsc{GovSim}-style: each action carries the
agent's own words and those words are what neighbours see. The second half is
forced rather than chosen---intervention tone is a manipulated factor and
language is its only carrier. A free group chat would make the intervener's tone
one voice among many and let citizens coordinate a joint response to it, so the
tone contrast would estimate the tone plus whatever the group negotiated about
it, which is not the estimand.

\subsection{From reply trees to visibility graphs}
\label{sec:visibility}

Agent adjacency is drawn from \textsc{Pheme} rumour threads. Each thread ships a
\texttt{structure.json} recording \emph{who replied to whom}; the environment
needs \emph{who can see whom}. These are not the same graph, and the step
between them is a modelling decision.

Taking the reply tree as the visibility graph and rooting the seven-node
neighbourhood at the cascade source produces a broadcast star most of the time,
because real \textsc{Pheme} cascades are broad and shallow (mean branching
$4.46$, mean depth $3.59$): a seven-node neighbourhood of the root is usually
the root plus six leaves. In a star rooted at the disseminator every citizen's
only neighbour is the disseminator, so no local majority can form among the five
citizens and \emph{no citizen can ever observe another paying the correction
cost}---the demonstration channel the headline outcome measures.

Two obvious repairs each fail in a different direction. Making all replies to a
common parent mutually visible fixes the connectivity by destroying locality: in
a star every reply is a sibling of every other, the neighbourhood becomes the
complete graph, and ``local majority'' becomes global majority. Sampling
neighbourhoods away from the root leaves most of them with no citizen visibility
at all.

A real thread view is neither extreme: it is ranked and truncated, so a reader
sees the source plus the replies near their own, not all fifty. We model that
directly---each reply is linked to the $w$ siblings following it in display
order---which is one parameter. At $w=1$ no neighbourhood on the real release is
a star, none is complete, mean density is $0.50$, and no thread is discarded.
We emphasise that $w=1$ is a \emph{proposal}: its justification is that it
produces non-degenerate neighbourhoods, which is not the same as being validated
against how threads are actually displayed. \S\ref{sec:rq1} therefore reports
all three translations rather than adopting one and suppressing the others.

\subsection{Experimental implementation}

One episode is one claim $\times$ one topology $\times$ one condition $\times$
one seed, run for $T = 12$ rounds with simultaneous moves, a payoff ledger
settled each round and veracity resolved at the horizon. The design crosses
intervention timing (early vs.\ late) $\times$ tone (empathetic nudge vs.\
aggressive debunking) $\times$ citizen backbone, against a no-intervener
control, a true-claim placebo arm, two ablations, an $N = 50$ scale check and a
perturbation arm, for $1{,}400$ episodes. Claims are stratified by
\textsc{Liar} severity and de-identified; from \textsc{Pheme} only derived
topology is used, never tweet content. Episode seeds are derived from a digest
so that a run replays exactly; the payoff ledger is unit-tested against
hand-computed episodes, because a silent ledger bug would invalidate every
primary outcome.

The results in \S\ref{sec:rq1}--\S\ref{sec:sensitivity} are produced with an
\emph{analytic surrogate} citizen policy---a bounded-rational model with
declared parameters---rather than a language model. This is deliberate for those
questions: \S\ref{sec:rq1} measures what a topology translation does to a
measured effect, and \S\ref{sec:rq2} measures what an accounting convention does
to a welfare ranking. In both, a language model is a source of variance, not the
object of study. It \emph{is} the object of study in \S\ref{sec:rq3}, which is
why that section is empty.

%======================================================================
\section{Experimental Results}

We organise results by research question.

\begin{description}
\item[RQ1] Does the reply-tree-to-visibility translation change the conclusion
  a study of this kind reaches? (\S\ref{sec:rq1})
\item[RQ2] Is total payoff a usable welfare measure in a design with costly
  sanctioning? (\S\ref{sec:rq2})
\item[RQ3] Do LLM societies catalyse or crowd out endogenous peer correction,
  and does the intervener's tone decide which? (\S\ref{sec:rq3})
\item[RQ4] How does LLM enforcement behaviour compare with human participants in
  public goods games? (\S\ref{sec:rq4})
\item[RQ5] Does an established correction norm survive a defector?
  (\S\ref{sec:rq5})
\end{description}

\paragraph{Metrics.}
\textbf{FPR}, the false-claim propagation rate, is the fraction of citizen
actions that propagate the claim. \textbf{EPC}, the endogenous peer-correction
rate, is the fraction that challenge it---that is, that voluntarily absorb
$\kappa$. Both are read directly off the action log; no model judge sits on the
critical path of either. Welfare is decomposed rather than summed
(\S\ref{sec:rq2}). Effect sizes are Hedges-corrected $d$; intervals are
bootstrap; multiplicity is Holm-corrected within each outcome family.

%----------------------------------------------------------------------
\subsection{RQ1: The visibility translation is the effect size}
\label{sec:rq1}

Table~\ref{tab:visibility} runs the identical core grid---same claims, same
threads, same seeds, same policies---three times, varying only how the reply
tree becomes a visibility graph.

\begin{table}[t]
\centering
\begin{tabular}{lrrrrr}
\toprule
Visibility model & No citizen & Mean & EPC & EPC & ATE [95\% CI] \\
 & visibility & reach & control & treated & \\
\midrule
Reply tree, rooted \emph{(as commonly specified)} & 88.6\% & 1.08 & 0.046 & 0.084 & $+0.038$ [+0.026, +0.050] \\
Bounded attention, $w=1$ \emph{(ours)} & 0.0\% & 2.19 & 0.028 & 0.109 & $+0.081$ [+0.067, +0.095] \\
All siblings visible & 0.0\% & 4.81 & 0.008 & 0.103 & $+0.095$ [+0.079, +0.112] \\
\bottomrule
\end{tabular}

\caption{The same corpus, the same seeds, three translations of reply structure
into visibility. The intervention's measured effect on endogenous peer
correction varies by a factor of $2.5$ across choices that papers in this area
do not usually report making. ``No citizen visibility'' is the share of
neighbourhoods in which no citizen can observe any other citizen.}
\label{tab:visibility}
\end{table}

Three things follow. First, the default---reply edges as visibility edges---is
not a neutral baseline but the most degenerate option available: it strands
88.6\% of neighbourhoods without any citizen-to-citizen edge, so the mechanism
the study measures is structurally absent from almost all of its data. Second,
the measured intervention effect is monotone in visibility density and ranges
over a factor of $2.5$; a paper reporting any one of these numbers without
stating the translation is reporting a quantity the reader cannot interpret.
Third, none of the environment's own quality gates catch this. The calibration
gate, which compares simulated cascade shape and the false-versus-true diffusion
asymmetry \citep{vosoughi2018} against the source threads, passes
\emph{comfortably} under the degenerate translation---it tests cascade shape,
and shape is insensitive to whether the agents can see one another.

We therefore recommend that studies seeding agent societies from cascade data
report the translation explicitly and, where feasible, report the sensitivity of
their headline effect to it, as here.

%----------------------------------------------------------------------
\subsection{RQ2: Summing a payoff function is not a welfare measure}
\label{sec:rq2}

\begin{table}[t]
\centering
\begin{tabular}{lrrl}
\toprule
Measure & No intervener & Intervener & Ranks intervention \\
\midrule
False-claim propagation rate (FPR) & 0.784 & 0.680 & \textbf{above} control \\
Endogenous peer correction (EPC) & 0.026 & 0.108 & \textbf{above} control \\
\midrule
$\sum \pi$ \emph{(total payoff)} & 51.283 & -74.127 & below control \\
$\sum \pi$ excluding sanction costs & 68.700 & 62.729 & below control \\
Conformity term $\beta C$ & 98.550 & 91.404 & below control \\
Veracity term $\gamma V$ & -29.850 & -28.675 & \textbf{above} control \\
\bottomrule
\end{tabular}

\caption{Control versus treated arms on the primary outcomes and on three
candidate welfare measures. The intervention cuts propagation and nearly
quadruples peer correction, and total payoff ranks it $122$ points below doing
nothing.}
\label{tab:welfare}
\end{table}

Table~\ref{tab:welfare} shows a measure that is mechanically anti-correlated
with its treatment. Reputation damage $\delta D$ is deadweight in the aggregate,
and generating the challenges that cause it is precisely the intervener's
function, so $\sum\pi$ is determined before any agent acts.

The obvious repair---report welfare with and without sanction costs, as is
standard in experimental public goods work \citep{fehr2002}---is necessary and
not sufficient. Excluding sanctions, the control arm still leads, because the
conformity term $\beta C$ dominates the remainder and conformity is maximised by
the arm in which nobody disagrees. We implemented three accounting conventions
for the ledger and all three preserve the ordering. The only component that
tracks the good being provided is the veracity term $\gamma V$, which is the
component we report as epistemic welfare.

This is not specific to \textsc{NudgeSim}. Any design in this line that reports
$\sum\pi$ as welfare, in a game where sanctioning is costly and the treatment
produces sanctioning, inherits the same inversion.

%----------------------------------------------------------------------
\subsection{RQ3: Benchmarking LLM agents}
\label{sec:rq3}

\pending{This section reports the model roster of \S\ref{sec:models} on FPR,
EPC, the welfare decomposition of Table~\ref{tab:welfare}, repair rate and
degeneracy statistics, one row per backbone, plus the timing $\times$ tone
interaction per backbone. It is the section for which the environment exists.
The grid is specified, the pipeline is validated (\S\ref{sec:sensitivity}), and
the run has not been executed; the measured cost of executing it is
\S\ref{sec:cost}. No number is reported here until it is.}

Table~\ref{tab:ate} shows the estimator that will populate it, applied to the
surrogate run: per-factor average treatment effects with bootstrap intervals,
following \citet{backmann2025moralsim}. Read as a property of the declared
data-generating process, not as a result about language models.

\begin{table}[t]
\centering
\begin{tabular}{llrr}
\toprule
Outcome & Factor & ATE [95\% CI] & Cohen's $d$ \\
\midrule
EPC & any intervention vs.\ no intervention & $+0.081$ [+0.067, +0.094]$^{*}$ & $+0.797$ \\
 & late entry vs.\ early entry & $-0.051$ [-0.070, -0.031]$^{*}$ & $-0.468$ \\
 & aggressive debunking vs.\ empathetic nudge & $-0.018$ [-0.038, +0.002] & $-0.166$ \\
\midrule
FPR & any intervention vs.\ no intervention & $-0.099$ [-0.119, -0.078]$^{*}$ & $-0.714$ \\
 & late entry vs.\ early entry & $+0.080$ [+0.054, +0.105]$^{*}$ & $+0.559$ \\
 & aggressive debunking vs.\ empathetic nudge & $+0.007$ [-0.020, +0.033] & $+0.047$ \\
\bottomrule
\end{tabular}

\caption{Per-factor average treatment effects with bootstrap 95\% intervals,
computed on the analytic surrogate. $^{*}$ significant after Holm correction
within the outcome family. \textbf{These are properties of the declared
data-generating process, not results about language models.}}
\label{tab:ate}
\end{table}

%----------------------------------------------------------------------
\subsection{RQ4: Comparison with human behaviour}
\label{sec:rq4}

Following \citet{piedrahita2025sanctsim}, results are read against human
participants rather than only against each other. We take the human row from
\citet{guerek2006} on two quantities both designs define: take-up of costly
enforcement, and sanctioning relative to affirmation. Absolute levels are not
comparable between a token-contribution game and a misinformation game, so the
comparison licenses an ordering and nothing stronger.

\begin{table}[t]
\centering
\begin{tabular}{llr}
\toprule
Measure & Population & Value \\
\midrule
Take-up of costly enforcement & Human \citep{guerek2006} & \textbf{0.929} \\
 & Analytic surrogate & 0.167 \\
\midrule
Sanctions per affirmation & Human \citep{guerek2006} & \textbf{1.660} \\
 & Analytic surrogate & 0.076 \\
\bottomrule
\end{tabular}

\caption{Human enforcement behaviour \citep{guerek2006} against the
\textsc{NudgeSim} control arm. The surrogate society barely enforces at all,
which is the second-order free-rider problem the environment is built to study,
measured rather than assumed.}
\label{tab:human}
\end{table}

\pending{The model rows of Table~\ref{tab:human}. \citet{piedrahita2025sanctsim}
find that LLMs invert the human pattern---rewarding where humans punish. Whether
the same inversion appears when the sanction is a public correction rather than
a token transfer is a direct test of whether that finding is about sanctioning
or about the specific game it was measured in.}

%----------------------------------------------------------------------
\subsection{RQ5: Norm robustness under a defector}
\label{sec:rq5}

Following \citet{piatti2024govsim}, one citizen defects at the half-way round,
either amplifying the false claim every round or refusing to pay the correction
cost thereafter. Because the society is sized by its topology neighbourhood, a
seat changing behaviour is the closest faithful analogue of \textsc{GovSim}'s
inserted newcomer. All outcomes exclude the perturbation seat in both arms, so a
fall means peers stood down rather than that the defector's own silence was
counted.

On the surrogate run, peers do not change their correction rate under either
defector. Read against Table~\ref{tab:human} that is coherent rather than
surprising: a society with almost no enforcement norm has none to defend.

\pending{The same contrast on LLM societies, which is where it carries
information: \textsc{GovSim} reports substantial degradation under perturbation,
and whether a correction norm is more or less robust than a harvesting norm is
the question this arm exists to answer.}

%----------------------------------------------------------------------
\subsection{Design sensitivity and pipeline validation}
\label{sec:sensitivity}

An analysis that has never been shown to detect an effect it was given is not
evidence about an effect it fails to find. We therefore run the pipeline against
three declared data-generating processes and require it to recover them in
order.

\begin{table}[t]
\centering
\begin{tabular}{lrrr}
\toprule
Tone$\to$EPC channel & seeds/cell & Cohen's $d$ & $n$/cell for 80\% power \\
\midrule
Declared & 30 & $-0.166$ & 572 \\
Inflated $\times 2.9$ & 30 & $-0.249$ & 254 \\
Set to zero & 30 & $-0.021$ & 35,322 \\
\midrule
Declared & 200 & $-0.061$ & 4,168 \\
Inflated $\times 2.9$ & 200 & $-0.203$ & 384 \\
Set to zero & 200 & $+0.037$ & 11,412 \\
\bottomrule
\end{tabular}

\caption{The analysis pipeline against declared data-generating processes.
Within each block the ordering is correct: inflating the channel enlarges the
effect, zeroing it leaves nothing.}
\label{tab:validation}
\end{table}

\begin{table}[t]
\centering
\begin{tabular}{lrrc}
\toprule
Contrast & Cohen's $d$ & $n$/cell for 80\% power & Powered at 30? \\
\midrule
Intervention $\to$ EPC & $+0.812$ & 25 & \checkmark \\
Intervention $\to$ FPR & $-0.748$ & 29 & \checkmark \\
H1 \quad timing $\to$ FPR & $-0.646$ & 39 & -- \\
H3 \quad tone $\to$ EPC \emph{(headline)} & $-0.148$ & 718 & -- \\
H2 \quad tone $\to$ FPR & $+0.028$ & 20,480 & -- \\
\bottomrule
\end{tabular}

\caption{Effect sizes and the per-cell sample each contrast needs for 80\%
power, measured on the surrogate grid. Only the intervention contrast is
powered at the 30 replications a conventional design allocates.}
\label{tab:power}
\end{table}

Table~\ref{tab:validation} also carries a caution. The declared effect is not
stable across claim draws---$-0.166$ at 30 replications per cell and $-0.061$ at
200---because a larger cell draws a different set of \textsc{Liar} claims rather
than more replications of the same ones. The tone contrast is under-powered in
the claim dimension as well as the seed dimension, and no single point estimate
should be read as its effect size.

Table~\ref{tab:power} is the reason the tone hypothesis is declared exploratory
before the grid runs rather than reinterpreted afterwards. It is also a general
caution for this literature: the intervention contrast is comfortable at 30
replications and the tone contrast needs nineteen times that, and nothing in the
design announces which is which.

\paragraph{A reproducibility failure worth naming.}
Episode seeds were originally derived from a hash of a string identifier.
Python salts string hashing per process, so two runs of the identical grid
agreed on every configuration and disagreed on $474$ of $600$ core-arm outcomes
while the run manifest recorded the seeds as pinned. Because each arm of
Table~\ref{tab:validation} ran in its own process, the declared, inflated and
zeroed processes were being compared across \emph{different} draws of claims and
topologies, and the inflated arm appeared to shrink the effect it inflates---an
ordering no parameterisation can produce, and therefore a bug rather than a
result. The general lesson is that a property a manifest asserts should be
enforced by a test, because a false assertion in a manifest is worse than no
assertion.

%----------------------------------------------------------------------
\subsection{What the benchmark costs}
\label{sec:cost}
\label{sec:models}

\begin{table}[t]
\centering
\begin{tabular}{lr}
\toprule
Quantity & Measured \\
\midrule
Model calls per episode & 60 \\
Mean input tokens per call & 616 \\
Calls per model over the grid & 23,250 \\
Duplicate prompts in 2,400 calls & \textbf{0\%} \\
Share of prompt that is a reusable prefix & 50.9\% \\
\midrule
Programme cost, recommended portfolio & \$2,055 \\
Programme cost, thinking on every arm & \$17,744 \\
\bottomrule
\end{tabular}

\caption{Metered from the prompts the environment assembles, priced against
published rates. Response caching---the lever such studies usually
name---returns nothing here, because every prompt carries a distinct claim,
agent and history.}
\label{tab:cost}
\end{table}

\paragraph{Model roster.}
The backbone factor has four levels: a frontier model with extended thinking
enabled, the \emph{same weights} with it disabled, a mid-tier model and a small
one. The paired frontier rows are the design's reason for existing:
\citet{piedrahita2025sanctsim} report their largest split between reasoning and
traditional models, but had to draw that contrast across different model
families, which confounds reasoning mode with everything else that differs
between them. Holding the weights fixed does not.

\paragraph{Cost.} Table~\ref{tab:cost} reports what that costs. Reasoning tokens
bill as output, so the thinking arm dominates: scheduling it for the main grid
and one high-power replication rather than for every calibration, validation and
re-run takes the programme from \$16{,}038 to \$1{,}866. That scheduling
decision is worth more than every caching lever combined. Cross-vendor breadth
comparable to \citet{backmann2025moralsim} or \citet{tewolde2026coopeval} is an
access question rather than an engineering one: a backbone is one line of
configuration, and the game, metrics and analysis layers are indifferent to
which model, or whether any model, supplies the citizen policies.

%======================================================================
\section{Limitations and Future Work}

\textbf{The model benchmark is not run.} Every claim in this paper is either
about the environment's design or about a declared analytic process. Nothing
here is evidence about language models, and the sections that would be are
marked and empty.

\textbf{The visibility model is a proposal, not a validated one.} Bounded
attention at $w=1$ is justified by producing non-degenerate neighbourhoods,
which is close to circular. What we establish is the sensitivity in
Table~\ref{tab:visibility}; what we do not establish is which translation is
correct. Validating $w$ against how threads are actually displayed and read is
the obvious next step.

\textbf{Seven agents.} \textsc{NudgeSim} is a controlled mechanism study with an
explicit payoff benchmark; an $N=50$ scale check provides directional
corroboration only, and scale dependence would be reported as a finding if it
appears.

\textbf{No negotiation channel.} The decision/observation split
(\S\ref{sec:env}) buys clean identification of the tone manipulation at the cost
of everything \textsc{GovSim} learns from group deliberation. Making the
institution endogenous---citizens voting on whether a corrector should exist and
who pays for it---is the natural follow-on and is where a deliberation channel
belongs.

%======================================================================
\section{Conclusion}

Casting counter-misinformation as a public goods game with costly peer
sanctioning turns ``did the intervention reduce sharing'' into a sharper
question: did it make the society more or less willing to pay for its own
correction. Building the environment to ask that question surfaced two problems
that are not specific to it. Turning reply structure into agent visibility is an
unreported researcher degree of freedom worth a factor of two and a half in the
headline effect, and the most obvious translation is the most degenerate one.
Summing a payoff function is not a welfare measure when the treatment's function
is to generate the costs being summed. Both apply to any study in this line, and
both are cheap to check before a compute budget is spent rather than after.

\bibliographystyle{plainnat}
\bibliography{refs}

\end{document}
